\newcommand{\wineIntro}{But we also remember all those of our generation who are not yet free.}

\newcommand{\wineSummaryThisYear}{}

\newcommand{\wineFirstCup}{
    \newReading The first cup represents physical freedom, the most basic freedom of all. Our ancestors could not accept the covenant of Torah while they were in Egypt because, as slaves, they lacked the freedom to determine the course of their own lives.

    Not all Jews in our world are free to live as Jews. We dedicate the first cup of wine to all those who still seek their physical freedom. As Jews, it is our responsibility to help them; this seder night we rededicate ourselves to that sacred purpose. May we soon share our freedom and our joy with all our brethren.
}

\newcommand{\wineSecondCup}{
    \newReading The second cup of wine symbolizes intellectual freedom\textemdash{}freedom of the mind. Closed minds lead to misunderstanding and human suffering. We need to open our minds to new ideas and try to understand the ideas and beliefs of others. Knowledge and understanding will lead to greater freedom for all in our world.
}

\newcommand{\wineThirdCup}{
    \newReading The third cup of wine symbolizes spiritual freedom. Our people have known the need for spiritual resistance in many ages. Even in the worst of circumstances, we have maintained our dignity. Even in the concentration camps, many Jews scrupulously observed \textit{Halakhah} (Jewish law) in defiance of the oppression they suffered at the hands of the Nazis. 
    
    \newReading One such Jew addressed this question to Rabbi Ephraim Oshry: ``Should a Jew, having to perform forced labor for the Nazis, continue to recite the benediction in the morning prayers: We praise You, Adonai our \god{}, ruler of the universe, who has not made me a slave?'' Rabbi Oshry responded: ``Heaven forbid that we should give up reciting this blessing. On the contrary, now of all times we are obliged to say this blessing so that our adversaries and tormentors realize that, although we are in their power to do with us as their wicked machinations devise, we nonetheless perceive ourselves not as slaves, but as free people, prisoners for the time being, whose liberation will soon come.''
    
    With this third cup of wine, let us seek the spiritual freedom that generations before us sacrificed to maintain. Let us open our hearts and minds to experience \god{} in our own lives.
}

\newcommand{\wineFourthCup}{
    \newReading As our seder draws near to the end, we take up our cups one last time. The redemption is not yet complete. Not everyone in our world is yet free. This fourth cup reminds us of our responsibility to be \god{}'s partners in bringing freedom to those enslaved, peace to those at war, food to those who hunger. This is our purpose as Jews. May we live to fulfill it.
}